% Subsection 5.1.2: 数据集与基线

\subsubsection{数据集口径}
训练集与评测集分别来自第三章顺序划分产出的 \texttt{data/processed/KodCode\_\allowbreak train.jsonl}（约 9000 条）与 \texttt{KodCode\_eval.jsonl}（约 1000 条）。两份文件在全部 10 组实验中共享，保证不同方法面对同一批题目，具体字段与清洗策略见 3.1--3.4 节。评测集样本均满足「单测可判定且参考解答通过」硬约束，避免无效样本拉低指标下限。

\subsubsection{基线与对照组}
以 H1（Base + Direct）作为共同基线，与之对照的 9 组覆盖训练侧与推理侧的完整析因：
\begin{itemize}
    \item \textbf{训练侧对照}：H4 / H7 / H9 分别引入 \texttt{LoRA\_code} / \texttt{LoRA\_cot\_zs} / \texttt{LoRA\_cot\_fs($k$)}，推理侧同为 Direct；
    \item \textbf{推理侧对照}：H2 / H3 在 Base 上叠加 CoT-ZS / CoT-FS($k$)；
    \item \textbf{交互对照}：H5 / H6 / H8 / H10 为训练 $\times$ 推理的联合配置，供第 5.3 节估计交互项。
\end{itemize}
全部配置与评测脚本入口参见 4.1 节表~\ref{tab:ch4-matrix}，本章直接按其口径取数不再重述。

\subsubsection{固定变量清单}
为让 10 组指标差异可归因到方法本身，以下变量在全部实验中固定不变：
\begin{itemize}
    \item \textbf{基座权重}、\textbf{tokenizer}、\textbf{chat template}、\textbf{数值精度}；
    \item \textbf{评测集}：同一份 \texttt{KodCode\_eval.jsonl}，不做抽样、不做重排；
    \item \textbf{采样配置}：每题采样数、温度与生成上限均固定；
    \item \textbf{执行环境}：同一子进程执行器与执行超时阈值；
    \item \textbf{few-shot 检索}：强标签集合、评分函数与评测侧 seed 派生方式。
\end{itemize}
其余所有变量——训练侧适配器种类、推理侧 prompt 结构、few-shot 示例数 $k$——均严格落入 4.1 节矩阵给定的取值，不做 grid 扩展。